\section{Discussion}

The central result of this study is implementation fidelity, not algorithmic novelty. After the final retained fixes, \tool{} reproduced MATLAB \blinker{} exactly across two public EEG resources that differ in experimental context, recording provenance and validation scale. That matters operationally: researchers can now move a mature ocular-index workflow into Python-native MNE pipelines without accepting silent changes in blink-event tables as the cost of migration.\supercite{gramfort2013,appelhoff2019}

The work also shows why a faithful port must be validated at the level of full-dataset reruns rather than spot checks. The remaining mismatches were not broad failures of the detector, but narrow differences in terminal-edge geometry and boundary-local floating-point behavior. Those are precisely the classes of software discrepancies most likely to survive superficial regression testing while still altering downstream analyses.

\paragraph{Self-critique and limitations.}
The strongest claim supported here is exact parity with the MATLAB reference, not superiority over alternative blink detectors and not independent confirmation that every blink boundary is physiologically optimal. First, the reference standard was MATLAB \blinker{} itself rather than a newly curated human-annotation gold standard, so any systematic bias present in the legacy implementation is inherited by this validation. Second, the evaluation spans only two public resources and the Murat benchmark used locally staged EDF/FIF-compatible derivatives rather than every original raw distribution format. Third, the final SEED result required a boundary-local numerical rescue in the validation harness after a floating-point mismatch was isolated near the recording end; this should be read as careful software replication, not as evidence that one globally tuned parameter set is universally sufficient. Fourth, we did not benchmark runtime, memory use, or the downstream stability of higher-order ocular indices beyond event boundaries.

Those limitations define the next steps. Independent human-annotation benchmarking, broader validation across additional device families and patient cohorts, direct assessment of downstream ocular-index reproducibility, and public archival release of a versioned software snapshot would all strengthen the package further. Even with those caveats, the present result is useful: a previously MATLAB-bound workflow can now be executed in Python with auditable, recording-level evidence that the migrated implementation preserves the legacy event tables exactly on the tested public datasets.
